\documentclass[11pt,a4paper]{article}
\usepackage{amsmath,amssymb,amsthm}
\usepackage{geometry}
\usepackage{hyperref}
\usepackage{cascade-xref}        % cross-paper \xref machinery (incl. xr-hyper)
\usepackage{booktabs}
\usepackage{graphicx}
\xrhyperdoc{part2}{cascade-series-part2}
\geometry{margin=2.5cm}

\newtheorem{theorem}{Theorem}[section]
\newtheorem{lemma}[theorem]{Lemma}
\newtheorem{corollary}[theorem]{Corollary}
\newtheorem{definition}[theorem]{Definition}
\newtheorem{principle}[theorem]{Principle}
\newtheorem{remark}[theorem]{Remark}

\title{The Cascade Series --- Prelude\\[0.5em]
\large Why Nothing Has Structure}
\author{RTAC}
\date{March 2026}

\begin{document}
\setlength{\emergencystretch}{3em}
\maketitle
\phantomsection
\label{paperdoc}

\begin{center}
\fbox{\parbox{0.9\textwidth}{\textbf{Status.} Foundational paper of the
cascade series. The prelude operates on finite-dimensional Euclidean
spaces only --- no infinite-dimensional Hilbert space is built or
required. The cascade's foundational footprint is pure
$\mathrm{RCA}_0$: natural numbers, computable functions on them, and
computable reals (Cauchy sequences with computable convergence rate).
Every analytic-looking quantity the cascade names --- $\pi$,
$\zeta(2)$, $V_d$, $\Omega_{d-1}$, the slicing-recurrence integral,
$\Gamma$ at half-integers, the $d\to\infty$ limits, the cascade
invariant $10^{-120}$ --- is a \emph{computable process}: a sequence of
rational approximations whose convergence to a target is
$\mathrm{RCA}_0$-expressible
($\forall\varepsilon\,\exists N$), without committing to a completed
real at the limit. The Monotone Convergence Theorem and stronger
analytic axioms are not imported as foundational commitments; they
serve as \emph{bookkeeping tools} in downstream papers
(Parts~0--VI), shorthand for the partial-computation form when
working through Gamma-function identities, integrals, and asymptotic
expansions. At any finite resolution depth (e.g., the universe's
current age in Planck ticks) every cascade prediction has a specific
rational value; ``exact identities'' are statements about shared
convergence targets, not about completed reals. Beyond what
counting on $\mathbb{N}$ itself requires --- the arithmetic content
of $\mathrm{RCA}_0$ that cumulation produces --- the cascade imports
no further axiom. The series'
load-bearing hypothesis --- that $B^\infty$ descended to four
dimensions is indistinguishable from our universe --- is stated
independently in the cover sheet and does not depend on the prelude's
chain.}}
\end{center}

\begin{abstract}
A theory of everything cannot have inputs: every input demands an
explanation of its origin, and that explanation is either circular or
requires a deeper theory. The only starting point that requires no
explanation is nothing. This prelude derives the cascade series'
starting structure --- the dimension tower
$\{(\mathbb{R}^d, B^d)\}_{d\in\mathbb{N}}$ of finite-dimensional
Euclidean spaces with asymptotic upper edge $B^\infty$ --- from a
single pre-mathematical input, the binary distinction $0\neq 1$. The
chain runs through cumulative distinguishability (binary, symmetric,
local), the countable cardinality tower forced by parameter economy
and minimal strength, the Euler product on $\mathbb{N}$ that forces
$\pi$ as a convergence target from the integer structure alone, and
the vector-space realisation $\mathbb{R}^d$ forced by Hurwitz's
theorem under a minimum-strength clause of austerity; the inner
product on $\mathbb{R}^d$ shares its angular convergence target with
the same prime-Euler $\pi$, and orthogonality between distinguishable
states is recovered as the extremal $\pi/2$ angular case forced by
(B), (S), (L). Every step is forced by an austerity clause or by a
classical theorem on integers and analysis. The cascade's foundational
footprint is pure $\mathrm{RCA}_0$: $\pi$, the slicing-recurrence
integral, the $d\to\infty$ limits, the cascade invariant $10^{-120}$,
and the rest of the analytic content are computable processes whose
convergence to a shared target is $\mathrm{RCA}_0$-expressible without
invoking a completed limit. The Monotone Convergence Theorem appears
only as a bookkeeping tool in downstream papers, not as a foundational
commitment. The starting structure carries a slicing recurrence with
unique constant $\sqrt{\pi}$, four distinguished dimensions
($5, 7, 19, 217$), and a cascade invariant of $10^{-120}$ ---
consequences of $0\neq 1$, not impositions on it. The series that
follows develops what these consequences imply about the geometry of
the world.
\end{abstract}

% ============================================================
\section{The Problem with Starting Points}\label{sec:starting-points}
% ============================================================

Every theory of everything proposed to date begins with something:
a string, a spin network, a set of fields, a symmetry group, a
Lagrangian, a wave function, a set of rules. Each starting point
carries structure, and that structure demands explanation.

Why a string rather than a membrane? Why $\mathrm{SU}(3)\times
\mathrm{SU}(2)\times\mathrm{U}(1)$ rather than $\mathrm{SU}(5)$?
Why ten dimensions rather than eleven? Every choice is a free input.
Every free input is a knob that could have been set differently. Every
knob that could have been set differently invites the question: what
set it?

The options are limited. Either the starting point is explained by
something deeper---in which case it was not the starting point---or it
is accepted as given, and the theory is incomplete. A theory of
everything that begins with something is, by construction, not a
theory of everything. It is a theory of everything \emph{given} its
starting point.

The alternative is to begin with nothing.

% ============================================================
\section{What ``Nothing'' Means}\label{sec:what-nothing-means}
% ============================================================

``Nothing'' in physics is usually taken to mean empty space, or the
quantum vacuum, or the absence of matter. These are not nothing. Empty
space has dimension, signature, topology. The quantum vacuum has
energy, fluctuations, fields. The absence of matter presupposes a
stage on which matter is absent.

Mathematical nothing is more austere: no space, no time, no
dimension, no content. But pure nothing---the absence of all
structure---cannot even be discussed, because discussion requires
the distinction between a statement and its negation.

This is the key observation. The minimal structure required for any
formal system to exist is not a space, not a symmetry, not a
dynamical law. It is the distinction between true and false.

\begin{definition}[Pre-mathematical axiom]\label{def:axiom}
$0 \neq 1$.
\end{definition}

This is not an axiom in the usual sense. It is not chosen from a
menu of alternatives. It is the precondition for axioms, theorems,
and proofs to be meaningful. A system that cannot distinguish $0$
from $1$ cannot state any proposition, cannot derive any consequence,
cannot describe any world. Definition~\ref{def:axiom} is the cost of
admission for coherent thought.

\begin{principle}[Austerity, derived from Definition~\ref{def:axiom}]\label{princ:austerity}
Under the commitment that Definition~\ref{def:axiom} ($0\neq 1$) is the
sole pre-mathematical input of the cascade series
(Section~\ref{sec:starting-points}), any construction introducing content
beyond what Definition~\ref{def:axiom} forces is inadmissible. This
excludes:
\begin{itemize}
\item[(i)] \emph{Parameter economy.} No free parameter is introduced ---
no unexplained numerical, categorical, or symmetry-group input.
\item[(ii)] \emph{Minimal strength.} No assumption strictly stronger than
Definition~\ref{def:axiom} is invoked where a weaker sufficient assumption
exists. Strength is measured in the sense of reverse mathematics over
the base theory $\mathrm{RCA}_0$
\cite{friedman1975,simpson2009}: an assumption $A$ is strictly
stronger than $B$ if $\mathrm{RCA}_0 + A \vdash B$ but
$\mathrm{RCA}_0 + B \nvdash A$.
\item[(iii)] \emph{Unobservable quotient.} Any continuous symmetry whose
action is unobservable from within the structure is quotiented out.
\end{itemize}
\end{principle}

\begin{proof}[Derivation from Definition~\ref{def:axiom}]
Each clause identifies a distinct way a construction can violate the
sole-input commitment of Definition~\ref{def:axiom}.

\emph{Clause (i).} A free parameter is a second explicit input alongside
Definition~\ref{def:axiom}. Admitting it contradicts the sole-input
status directly.

\emph{Clause (ii).} If a construction uses assumption $A$ where the
weaker $B$ suffices (with $A\Rightarrow B$ strictly), the residue
$A\setminus B$ is additional content not forced by
Definition~\ref{def:axiom}. That residue is a second input --- the same
class of violation as clause~(i).

\emph{Clause (iii).} A continuous symmetry whose action is unobservable
from within the structure carries structural content with no derivation
consequence. Retaining it imports content not forced by
Definition~\ref{def:axiom} --- the same class of violation.

\emph{Exhaustiveness of (i)--(iii) (formal-extension trichotomy).} A
construction $C$ violates sole-input iff $C$ is derivable in some
formal system $\Sigma_C$ properly extending $\Sigma_0 = (\text{logic})
+ \{0\neq 1\}$. Standard proof theory characterises a formal system by
three independent axes --- its non-logical axioms, its underlying
logic, and its language/signature. Any extension of $\Sigma_0$ acts on
at least one of these axes:
\begin{itemize}
\item Extending the axioms imports content as a fresh assumption
(numerical, categorical, or symmetry-group input). This is clause~(i).
\item Extending the logic imports content as inferential strength
beyond what Definition~\ref{def:axiom} requires (e.g.\ adopting
classical reasoning, the axiom of choice, or higher-order quantification
where the cascade's chain demands only weaker resources); the residue
$A\setminus B$ between the adopted logic $A$ and the strictly weaker
sufficient $B$ encodes the imported content. This is clause~(ii).
\item Extending the signature or background imports content as ambient
structure with no cascade-internal observable trace --- a continuous
symmetry whose action is invisible from within the cascade. This is
clause~(iii).
\end{itemize}
These three axes exhaust the formal anatomy of an extending system:
there is no fourth axis along which a formal system can be extended in
standard proof theory. Any candidate fourth clause therefore either
reduces to one of (i)--(iii) under the standard decomposition or fails
to identify content actually present in $\Sigma_C$'s extension of
$\Sigma_0$. Definitional extensions of $\Sigma_0$ are conservative
(they introduce notation but no new theorems) and therefore do not
qualify as content additions; they require no clause.

All three clauses track a single condition: admitting content not forced
by the sole input. Austerity is therefore the operational content of
Definition~\ref{def:axiom}'s status as sole pre-mathematical input, not
a separate axiom.
\end{proof}

\noindent\textbf{Status.} Principle~\ref{princ:austerity} is a
\emph{derivation rule} for cascading from $0\neq 1$, analogous to
modus ponens being the inference rule of propositional logic: it is
built into what ``sole pre-mathematical input'' means, rather than a
further axiom within the mathematical content. Each step of the prelude
chain that previously relied on an implicit parsimony choice is now
annotated with the clause of Principle~\ref{princ:austerity} it invokes.
The chain from $0\neq 1$ to the cascade's complete starting structure
is closed at every step: each step is forced by a clause of
Principle~\ref{princ:austerity} or by a classical theorem on integers
and analysis (Section~\ref{sec:the-chain}).

The question is: what does $0\neq 1$ \emph{imply}?

% ============================================================
\section{Distinction and Cumulation}\label{sec:distinction-and-cumulation}
% ============================================================

If $0$ and $1$ are distinguishable, they share no content. The
logical NOT operation, mapping $0\mapsto 1$ and $1\mapsto 0$, is an
involution exchanging two states with no common component. This
section unpacks what \emph{distinguishable} means structurally, at
the combinatorial level. The number-theoretic content that follows
from cumulation is the subject of
Section~\ref{sec:integer-structure}; the geometric content (vector
spaces, inner products, orthogonality at $\pi/2$) is derived in
Section~\ref{sec:numbers-to-geometry}, after $\pi$ has been forced from
the integer structure.

\subsection*{Three structural properties of $\neq$}

The relation $\neq$ has three structural properties that follow from
its meaning as a binary relation in standard logic, not from any
cascade-internal commitment:
\begin{itemize}
\item[(B)] \emph{Binary.} Distinguishability is all-or-nothing: two
states either are or are not distinguishable. There is no
``partially distinguishable'' --- $\neq$ does not admit gradations.
\item[(S)] \emph{Symmetric.} $a \neq b$ iff $b \neq a$.
\item[(L)] \emph{Local.} The fact that $a \neq b$ depends only on
$a$ and $b$, not on what other distinguishable elements exist. The
relation between $0$ and $1$ is the same whether or not $2, 3,
\ldots$ also exist.
\end{itemize}

(B), (S), and (L) are operational content of ``$\neq$ is a binary
relation in standard logic''; they require no commitment beyond
Definition~\ref{def:axiom}.

\subsection*{Cumulation}

Each new distinguishable state extends the cumulative pool of
pairwise distinctions: adding $s_{d+1}$ to a set of $d$
pairwise-distinguishable states adds $d$ new pairwise distinctions
(one with each existing state), each binary, symmetric, and local.
Iterating, the cumulation of $d$ pairwise-distinguishable states
forms a $d$-element discrete structure with full pairwise distinction
--- combinatorially the vertex set of the complete graph $K_d$, with
$\binom{d}{2}$ pairwise distinctions determined entirely by the
cardinality $d$. At this stage no algebraic combination is imported;
the ``combination'' of states is the bookkeeping of cumulative
distinction itself. No vector space, no inner product, no notion of
angle has yet been introduced.

\subsection*{What cumulation gives us}

Cumulation produces an integer count: the cardinality $d$ of the
distinguishable state set. Integers are not separately imported ---
they are what cumulation \emph{is}, the counting data of the discrete
structure. The cascade's downstream content (the geometric
realisation, the slicing recurrence, the cosmological constant) all
follow from this integer cardinality together with the computable
functions $\mathrm{RCA}_0$ supplies on it. Number theory enters in
Section~\ref{sec:integer-structure}; geometry in
Section~\ref{sec:numbers-to-geometry}. Cumulation by itself is purely
combinatorial.

% ============================================================
\section{The Integer Structure: Cardinality and \texorpdfstring{$\pi$}{pi}}\label{sec:integer-structure}
% ============================================================

Cumulation produces an integer count: the cardinality $d$ of
distinguishable states. The existence of one distinction does not
preclude a second; the cascade's distinguishable-state set grows
without bound, indexed by a positive integer at each stage of the
descent. This section establishes two things: (i)~the cardinal
indexing of the cascade's tower is forced to range over $\mathbb{N}$
by clauses (i) and (ii) of austerity; (ii)~the integer structure of
$\mathbb{N}$ forces $\pi$ as a convergence target via the Euler
product on primes --- a computable process whose partial values are
$\mathrm{RCA}_0$-expressible without any commitment to a completed
real at the limit. Both results are number-theoretic; neither imports
any geometric content.

\subsection*{Cardinality}

\begin{theorem}[Countable dimension tower]\label{thm:iterate}
The cascade's dimension parameter $d$ ranges over $\mathbb{N}$. The
choice of $\mathbb{R}^d$ as the continuous realisation medium at each
$d$ is derived in Section~\ref{sec:numbers-to-geometry}; this theorem
establishes the cardinal indexing $d\in\mathbb{N}$ alone.
\end{theorem}

\begin{proof}
\emph{Existence of arbitrarily many cardinalities.} For any finite
$d$, the $d$-element distinguishability record (the $K_d$-shaped
combinatorial structure of Section~\ref{sec:distinction-and-cumulation})
is expressible in $\mathrm{RCA}_0$ \cite{simpson2009} --- the
strength-measurement base of Principle~\ref{princ:austerity}
clause~(ii).

\emph{The tower has no upper bound (clause~(i)).} Suppose the tower
stops at some finite $n<\infty$. Then $n$ itself is a structural
input of the theory --- an unexplained integer specifying the
cascade's maximum cardinality. By Principle~\ref{princ:austerity}
clause~(i) (parameter economy), any such integer is rejected as a
free parameter. The only stopping points compatible with clause~(i)
are $n=0$ (which contradicts Definition~\ref{def:axiom}, requiring at
least two distinct states, hence cardinality at least~$1$) and
$n=\infty$. The tower therefore extends without bound. (``$n=\infty$''
here is not itself a chosen cardinality among alternatives: within
the $\mathrm{RCA}_0$ strength-measurement base, $\aleph_0$ is the
unique unbounded notion --- $\mathrm{RCA}_0$ has no machinery to
express higher cardinals --- so $n=\infty$ is the unbounded countable
index $\mathbb{N}$ the non-termination directly produces, not a
selection from a menu of infinities.)

\emph{Countable rather than uncountable indexing (clause~(ii)).}
Indexing the tower by an uncountable cardinal would require
expressive machinery beyond second-order arithmetic, strictly
stronger than Definition~\ref{def:axiom} in the reverse-mathematics
sense of clause~(ii) \cite{friedman1975,simpson2009}.
Principle~\ref{princ:austerity} clause~(ii) (minimal strength)
therefore excludes the uncountable case. The cardinality $d$ ranges
over $\mathbb{N}$.
\end{proof}

\subsection*{Number theory on \texorpdfstring{$\mathbb{N}$}{N}: \texorpdfstring{$\pi$}{pi} via the Euler product}

The cascade's downstream content --- the slicing-recurrence integral
$\int_{-1}^1 (1-x^2)^{d/2}\,dx$, sphere areas, ball volumes,
asymptotic behaviour as $d\to\infty$ --- has an analytic look but,
as developed below, lives entirely on computable functions of
$\mathbb{N}$. Each such quantity is a sequence of rational
approximations indexed by a resolution parameter $N$: partial sums,
partial products, Riemann sums, truncated improper integrals. The
$\mathrm{RCA}_0$-expressible content of the sequence is its
\emph{convergence behaviour} --- statements of the form ``for every
rational $\varepsilon>0$ there exists a computable $N$ such that the
$N$th approximation is within $\varepsilon$ of all subsequent
approximations.'' We do not commit to a completed real at the limit;
we commit only to the partial-computation form, which is
$\mathrm{RCA}_0$-internal.

By Principle~\ref{princ:austerity} clause~(ii) (minimal strength),
this is the austere reading: stronger axioms (the Monotone
Convergence Theorem and its kin) would \emph{also} suffice and would
let us name the targets as completed reals, but they import content
the cascade does not require. The cascade's foundational footprint
stays at pure $\mathrm{RCA}_0$; $\mathrm{MCT}$ and stronger
principles appear only as \emph{bookkeeping tools} in downstream
papers, shorthand for the partial-computation form when working
through Gamma identities and asymptotic expansions.

\begin{theorem}[The integer structure forces $\pi$ as a convergence target]\label{thm:pi-from-primes}
On the cumulation's integer structure $\mathbb{N}$, the Euler product
over primes defines a computable sequence of rational approximations
whose convergence to a shared target is $\mathrm{RCA}_0$-expressible.
Writing the partial sums and partial products
\[
\zeta_N(2) \;=\; \sum_{n=1}^{N} \frac{1}{n^2},
\qquad
\mathcal{E}_N \;=\; \prod_{\substack{p\,\mathrm{prime}\\p\leq N}}
\frac{1}{1 - p^{-2}},
\qquad
\pi_N^2 \;=\; 6\,\zeta_N(2),
\]
the cascade's $\pi$ is the convergence target of $\sqrt{\pi_N^2}$
under the $\mathrm{RCA}_0$-internal convergence statement
$\forall\varepsilon\,\exists N$ such that $|\pi_M - \pi_N|<\varepsilon$
for all $M\geq N$. The classical identities
$\zeta(2) = \prod_p (1 - p^{-2})^{-1}$ and $\zeta(2)=\pi^2/6$ are
shorthand for ``the three sequences $\zeta_N(2)$, $\mathcal{E}_N$,
and $\pi_N^2/6$ converge to the same target,'' a statement provable
in $\mathrm{RCA}_0$ from the term-by-term identity
$\zeta_N(2) = \mathcal{E}_N + O(1/N)$ (geometric-series expansion and
the fundamental theorem of arithmetic). The integer structure of
$\mathbb{N}$ alone, without any geometric input, forces this
convergence target.
\end{theorem}

\begin{proof}
The partial sums $\zeta_N(2)$ and the partial Euler products
$\mathcal{E}_N$ are computable functions of $N$ definable in
$\mathrm{RCA}_0$. Their term-by-term agreement follows from the
finite-product expansion of $\mathcal{E}_N$ as a sum over integers
whose prime factors are all $\leq N$, which agrees with
$\zeta_N(2)$ up to terms of order $1/N$ (a bound provable in
$\mathrm{RCA}_0$). The Basel-Euler identity ($\zeta_N(2)$ and
$\pi_N^2/6$ converging to the same target) reduces to the classical
1734/1737 proofs read at finite resolution: each Wallis, Fourier,
or sin-product proof produces partial-sum identities matching to
$O(1/N)$ at every $N$, and the Cauchy form of convergence
($\forall\varepsilon\,\exists N$) is then $\mathrm{RCA}_0$-internal.
Uniqueness of the target is elementary: any two computable sequences
sharing the same Cauchy convergence record determine the same target
up to $\mathrm{RCA}_0$-isomorphism.
\end{proof}

\begin{remark}[$\pi$ enters from number theory, not from geometry; as a process, not as a completed real]
Theorem~\ref{thm:pi-from-primes} forces $\pi$ as a convergence
target via the integer/prime structure of $\mathbb{N}$ within
pure $\mathrm{RCA}_0$. No vector space, no inner product, no angle,
no Euclidean structure has yet been introduced. The identification
of this Euler-product $\pi$ with the angular convergence target
($\pi/2$ for orthogonality, $2\pi$ for a full circle, $\pi$ for the
half-circle) is a downstream consequence
(Section~\ref{sec:numbers-to-geometry}): when geometry is derived,
its $\pi$-sequence is the same one, term-by-term. The cascade's
Euclidean geometry inherits $\pi$ from number theory rather than
supplying it. Importantly, $\pi$ here is not a static completed
real --- it is the cascade's open-ended computation, with the
universe's current age in Planck ticks corresponding to a finite
resolution depth $N$ at which the partial value $\pi_N$ is a
specific rational. The completed-real reading is a downstream
bookkeeping convenience, not a commitment of this prelude.
\end{remark}

\noindent\textbf{No infinite-dimensional ambient space.} Theorem~\ref{thm:iterate}
does not build a single infinite-dimensional Hilbert space as an
ambient object. The cascade is the dimension tower itself: the
sequence of $K_d$-shaped discrete structures (extended to
$\mathbb{R}^d$ in Section~\ref{sec:numbers-to-geometry}) indexed by
$d\in\mathbb{N}$, together with the relations between consecutive
levels (the slicing recurrence of
Section~\ref{sec:nothing-unstable}). The
``infinite-dimensional unit ball $B^\infty$'' is the asymptotic upper
edge of the tower, defined as the $d\to\infty$ limit of intrinsic
geometric quantities (volumes, sphere areas; see
Theorem~\ref{thm:nothing} below), not as a literal
infinite-dimensional ambient object. This means the cascade does not
require Cauchy completion to a single infinite-dimensional Hilbert
space, does not require any functional-analytic refinement (rigged,
nuclear, weak-topology), and does not require Sol\`er-style
convergence theorems to force its starting structure: the cardinal
indexing of the tower is forced by clauses~(i) and~(ii) of
Principle~\ref{princ:austerity} (Theorem~\ref{thm:iterate}); $\pi$
is forced by the Euler product on $\mathbb{N}$ within pure
$\mathrm{RCA}_0$ as a convergence target
(Theorem~\ref{thm:pi-from-primes}); the choice of $\mathbb{R}^d$ as
the geometric realisation is derived in
Section~\ref{sec:numbers-to-geometry}. The cascade's entire
foundational footprint stays at pure $\mathrm{RCA}_0$ for both the
cardinal tower's existence and the classical-analysis content
(integration, the Gamma function, $d\to\infty$ limits) developed in
subsequent sections --- each analytic identity is a convergence
statement about computable sequences, expressible in $\mathrm{RCA}_0$
without invoking the Monotone Convergence Theorem or any stronger
analytic principle. $\mathrm{MCT}$ and its kin appear only as
bookkeeping shorthand in Parts~0--VI.

% ============================================================
\section{From Numbers to Euclidean Geometry}\label{sec:numbers-to-geometry}
% ============================================================

Section~\ref{sec:integer-structure} forced the cardinality
$d\in\mathbb{N}$ from cumulation and clauses~(i)--(ii) of austerity,
and forced $\pi$ as a convergence target from the Euler product
within pure $\mathrm{RCA}_0$. Both results live on the integer
structure of $\mathbb{N}$: no vector space, no inner product, no
angle has yet been introduced. This section closes the gap from
that number-theoretic content to a Euclidean geometric realisation:
the vector-space ambient $\mathbb{R}^d$, its inner product, and the
extremal pairwise angle $\pi/2$ at which the original $d$
distinguishable states sit. Every step is forced by clause~(ii) of
austerity or by a classical theorem.

\subsection*{Why a continuous realisation}

$\mathrm{RCA}_0$ supplies the computable reals --- a countable
collection of computable Cauchy sequences with computable
convergence rate, sufficient to host every metric, vector-space, and
inner-product structure the cascade names. Cumulation
(Section~\ref{sec:distinction-and-cumulation}) is silent on
questions of magnitude --- it records whether states are
distinguishable, not \emph{how} --- so the cumulated states need
this computable-real container before any non-binary geometric
relation between them can be expressed. The question is what
\emph{shape} that container takes. All ``continuous'' content below
is read in the partial-computation form of
Section~\ref{sec:integer-structure}: each magnitude is a computable
sequence with an $\mathrm{RCA}_0$-expressible convergence record.

\subsection*{The vector-space realisation}

A geometric realisation that quantifies how much two states overlap
requires the realised states to live in a vector space supporting an
inner product. By Hurwitz's theorem, the associative normed division
algebras over $\mathbb{R}$ are exactly $\mathbb{R}, \mathbb{C},
\mathbb{H}$. By Principle~\ref{princ:austerity} clause~(ii)
(minimal strength), $\mathbb{R}$ is the minimum: starting over
$\mathbb{C}$ or $\mathbb{H}$ imports complex structure ($J^2 =
-\mathrm{Id}$) or quaternionic structure ($i, j, k$ with their
relations) as additional axiomatic content beyond what
Definition~\ref{def:axiom} and the cumulated integer structure of
$\mathbb{N}$ force. The vector-space realisation is therefore
$\mathbb{R}^d$ --- the cardinal $d$ supplied by
Theorem~\ref{thm:iterate}, the field supplied by Hurwitz +
clause~(ii).

The Euclidean inner product on $\mathbb{R}^d$ is the unique
positive-definite symmetric bilinear form up to isomorphism, given
a basis (elementary linear algebra); it is fixed once the field is
chosen, not separately committed. With the inner product comes its
geometric interpretation: the cosine of an angle. The angle between
unit vectors ranges over the convergence-target interval $[0, \pi]$,
with $0$ corresponding to identical states, $\pi/2$ to orthogonal,
$\pi$ to antipodal. The angular convergence target $\pi$ here is
the same computable sequence as the prime-Euler $\pi$ forced by
Theorem~\ref{thm:pi-from-primes}: the geometric definition produces
the same partial-approximation sequence (term-by-term, at every
resolution depth $N$) as the Euler-product definition, so the
geometric and number-theoretic $\pi$'s coincide as processes, not
just as completed values. No new analytic content is introduced.

\subsection*{The extremal angular case: pairwise orthogonality}

Within the $\mathbb{R}^d$ realisation, the $d$ cumulated states sit
at a forced pairwise angle.

\begin{theorem}[Pairwise orthogonality of cumulated states]\label{thm:orth}
Let the $d$ cumulated distinguishable states of
Section~\ref{sec:distinction-and-cumulation} be realised as
unit vectors in $\mathbb{R}^d$ (the realisation forced above). Their
pairwise inner products are $\langle e_i, e_j\rangle = \delta_{ij}$:
distinct cumulated states are orthogonal, and the realisation is the
standard basis of $\mathbb{R}^d$.
\end{theorem}

\begin{proof}
Property (B) (binary distinguishability) of $\neq$ admits no
gradation: between two cumulated states, the inner product takes an
extremal value on the scale $[-1, 1]$, not an intermediate one. The
available extrema are $+1$ (identical), $0$ (orthogonal), and $-1$
(antipodal): $+1$ is excluded because the states are distinct, $-1$
is excluded because it imposes a privileged pairing structure
(antipodal partners) that has no analogue in (B) ---
distinguishability is a relation between two states, not a coupling
that links one state to a unique opposite. The extremal value
compatible with all three of (B), (S), and (L) is $0$.

Property (L) (locality) requires that the pairwise relation depend
only on the pair, not on the cardinality of the state set: this
excludes the regular-simplex realisation in $\mathbb{R}^{d-1}$, whose
pairwise inner product $-1/(d{-}1)$ depends explicitly on $d$.
Property (S) (symmetry) is satisfied by any inner product. The
realisation forced by (B), (S), and (L) is therefore the standard
basis of $\mathbb{R}^d$ with $\langle e_i, e_j\rangle = \delta_{ij}$
and pairwise angle $\pi/2$ --- the extremal ``no overlap'' value of
the inner-product scale.
\end{proof}

\begin{remark}[$\pi/2$ inherited from $\pi$, not separately imported]
The pairwise angle $\pi/2$ is half of the angular extent $\pi$
forced by Theorem~\ref{thm:pi-from-primes}. No additional axiom is
invoked: the integer structure of $\mathbb{N}$ supplied $\pi$ via
the Euler product; Hurwitz + clause~(ii) supplied $\mathbb{R}^d$;
linear algebra supplied the inner product; (B), (S), (L) supplied
the extremal orthogonal value. The chain from number theory to
Euclidean geometry closes without circling back through any
geometric assumption.
\end{remark}

\medskip

\noindent The cascade's downstream $\sqrt{\pi}$ constant in the
slicing recurrence (Section~\ref{sec:nothing-unstable}) follows: the
orthogonal angle $\pi/2$ between axis and equator forces the
half-integer argument in $B(1/2,\cdot)$, giving $\Gamma(1/2) =
\sqrt{\pi}$. The chain --- $\pi$ forced via the Euler product on
$\mathbb{N}$ (Section~\ref{sec:integer-structure}), vector-space
realisation forced by Hurwitz + clause~(ii), inner product forced by
linear algebra, pairwise orthogonality forced by (B), (S), (L) ---
is enumerated together with the rest of the prelude's chain in
Section~\ref{sec:the-chain}.

% ============================================================
\section{Scale and Rotation Invariance: Collapse at Each Dimension}\label{sec:scale-and-rotation-invariance-collapse-a}
% ============================================================

Starting from nothing, there is no external reference against which to
measure absolute magnitude. There is no metre stick, no Planck mass,
no energy scale. The only quantities that can be defined without an
external reference are \emph{ratios}. At each $d$, the dilation
$x\mapsto\lambda x$ on $\mathbb{R}^d$ for $\lambda\in\mathbb{R}^{*}$
is unobservable from within the structure. By
Principle~\ref{princ:austerity} clause~(iii) (unobservable quotient),
the cascade structure at each $d$ is invariant under dilation.

The same applies to rotations. The orthogonal group $\mathrm{SO}(d)$
acts on the labeling of the orthonormal frame in $\mathbb{R}^d$, not
on the frame's existence. Cascade calculations downstream (the slicing
recurrence of Section~\ref{sec:nothing-unstable}, the four
distinguished dimensions of Part~0, the cascade invariant) are
$\mathrm{SO}(d)$-invariant: they depend on intrinsic geometric
quantities (volumes, sphere areas, Gamma function values), not on any
specific axis. By clause~(iii), $\mathrm{SO}(d)$ is also quotiented at
each $d$.

\noindent\textbf{The collapsed substrate at each $d$.} The combined
quotient $\mathbb{R}^d/(\mathbb{R}^*\times\mathrm{SO}(d))$ collapses to
a single point: $\mathrm{SO}(d)$ acts transitively on rays in
$\mathbb{R}^d$, so all rays are identified; $\mathbb{R}^*$ identifies
points along each ray. At every finite $d$, the projective space modulo
rotations is therefore a single point. This collapse is not a vacuous
result --- it is the cascade's substrate at each $d$. Structured
nothing has full rotational symmetry: no direction is privileged, no
basis is canonical, and the only data the substrate carries at level
$d$ is the bare fact of being a $d$-dimensional unit ball. The
dimension parameter $d$ survives both quotients (it is counting data,
not a direction), and at each $d$ the intrinsic geometric content of
the unit ball $B^d$ --- its volume $V_d$, its boundary sphere area
$\Omega_{d-1}$, the slicing-recurrence ratio $R(d)$, the decay rate
$p(d)$, the effective slicing radius $R_{\mathrm{eff}}(d)$ --- is
rotation-invariant and non-trivial.

\noindent\textbf{The cascade is the resolution of unresolved infinity.}
The cascade is not a static structure indexed by $d$. It is the
\emph{process} by which the unresolved infinity at the asymptotic
upper edge of the dimension tower is resolved into finite-dimensional
intrinsic content, dimension by dimension. The slicing recurrence
(Section~\ref{sec:nothing-unstable}) is the resolution mechanism:
each step relates the intrinsic content at $d+1$ to the intrinsic
content at $d$, extracting one dimension's worth of finite content
from the asymptotic limit. The dimension tower
$\{(\mathbb{R}^d, B^d)\}_{d\in\mathbb{N}}$ is the trace of this
resolution; the integer $d$ is the resolution stage. The universe
does not contain infinity --- the universe \emph{is} the result of
resolution applied to it, observed from a particular stage. This is
the unifying observation of the cascade series (cover sheet, ``one
infinity, not many''): every infinity problem in physics --- the
cosmological constant's $10^{-120}$, the black-hole horizon's
infinite time dilation, the ultraviolet floor at $\Omega_{217}$
(manifest temporally as the Big Bang in the universe's past) --- is a
local slice of the same asymptotic boundary, and every finite
structure observed is what that one infinity looks like partway
through resolution.

\noindent\textbf{Computational representative.} Although the substrate
at each $d$ is the collapsed quotient, calculations are performed on
$\mathbb{R}^d$ directly, with its standard inner product and unit ball
$B^d$. Cascade calculations on $\mathbb{R}^d$ are
$\mathrm{SO}(d)$-equivariant by construction --- they descend to the
collapsed substrate because their content (volumes, areas, dimension
counts, Gamma function identities) is invariant under any choice of
frame. Unit norm is the choice of canonical representative of each
ray; the choice of frame within $\mathbb{R}^d$ is a further gauge
fixing that washes out under $\mathrm{SO}(d)$ when the calculation
completes.

\noindent\textbf{Other candidate symmetries.}
\begin{itemize}
\item \emph{$\mathbb{C}^*$-phase rotation}. A phase action requires a
complex structure $J$ with $J^2=-\mathrm{Id}$. No such structure is
present at the Prelude stage: each $\mathbb{R}^d$ is constructed over
$\mathbb{R}$ (see Section~\ref{sec:what-not}, ``Real start, complex
amplitudes, quaternionic spacetime''). The forced precession
$\alpha=\pi/2$ of Paper~II induces $J$ downstream, at which point
whether $\mathbb{C}^*$ is observable or unobservable becomes a
separate question for Paper~II and beyond. At the Prelude stage there
is no $\mathbb{C}^*$ action to quotient.
\item \emph{Discrete scale invariance} (quotient by a specific ratio
$\lambda = r$ for some fixed $r$). Selecting a specific $r$ introduces
an unexplained numerical input, violating clause~(i) (parameter
economy). A continuous quotient has no such free parameter; only
discrete reductions introduce one.
\end{itemize}

\begin{definition}[Sphere representative and ball completion]\label{def:ball}
The unit sphere $S^{d-1}=\partial B^d$ is the canonical 2-fold
representative of the projective space
$\mathbb{R}\mathbb{P}^{d-1}=(\mathbb{R}^d\setminus\{0\})/\mathbb{R}^*$:
each ray through the origin meets $S^{d-1}$ in exactly one antipodal
pair. The unit ball $B^d = \{x\in\mathbb{R}^d : |x|\leq 1\}$ is the
convex completion: each ray meets $B^d\setminus\{0\}$ in a half-open
radial segment rather than at a single point, so the ball is not
itself a section of $\mathbb{R}\mathbb{P}^{d-1}$ but contains
$S^{d-1}$ as its boundary. The ball organises the slicing recurrence
(Remark below); the sphere is where the cascade's content lives. Every
subsequent theorem of the cascade is projective-invariant and
rotation-invariant --- it depends only on dimension counts and intrinsic
geometric ratios, not on the chosen representative or its completion
or any specific frame --- so the joint choice is computational, not
ontological.
\end{definition}

\begin{remark}[Why the ball rather than the sphere]
We work with the ball $B^d$ rather than the sphere $S^{d-1}$ because
the ball is the natural site for the slicing recurrence: slicing
$B^{d+1}$ perpendicular to any axis at height $x\in[-1,1]$ yields a
ball $B^d$ of radius $\sqrt{1-x^2}$, whose boundary is a sphere
$S^{d-1}$. This is the slicing recurrence of
Section~\ref{sec:nothing-unstable}, and it is the mechanism by which
the Gamma function enters. The sphere $S^d$ does not carry an
analogous native recurrence to lower-dimensional spheres; the ball is
the minimal container in which the recurrence is expressible without
auxiliary structure. Once the recurrence is in place, boundary
dominance ($\Omega_{d-1}/V_d = d$) shows that the spheres carry
essentially all the content --- so physics lives on the sphere, but
the recurrence is organised by the ball.
\end{remark}

\begin{corollary}[The starting structure]\label{cor:start}
The mathematical structure corresponding to ``nothing with the
capacity for distinction'' is the dimension tower
$\{(\mathbb{R}^d, B^d)\}_{d\in\mathbb{N}}$ of finite-dimensional
Euclidean spaces and their unit balls, with each level collapsed to
a single substrate point under
$\mathbb{R}^*\times\mathrm{SO}(d)$ but carrying non-trivial intrinsic
geometric content (volumes, areas, recurrence ratios). The asymptotic
upper edge --- the $d\to\infty$ limit, denoted $B^\infty$ --- is
structured nothing (Theorem~\ref{thm:nothing} below). The cascade is
the trajectory of intrinsic content unfolded across this tower via
the slicing recurrence (Section~\ref{sec:nothing-unstable}).
\end{corollary}

% ============================================================
\section{Nothing Is Empty}\label{sec:nothing-is-empty}
% ============================================================

The asymptotic upper edge of the dimension tower --- denoted
$B^\infty$, the convergence target of the intrinsic geometric
quantities of $B^d$ as $d$ grows --- is empty in every measurable
sense.

\begin{theorem}[Geometric nothing at the asymptotic upper edge]\label{thm:nothing}
The volume and sphere-area sequences $V_d$ and $\Omega_{d-1}$ are
computable functions of $d$ and converge to zero under the
$\mathrm{RCA}_0$-internal Cauchy condition: for any rational
$\varepsilon>0$ there exists a computable $N$ such that $V_d <
\varepsilon$ and $\Omega_{d-1} < \varepsilon$ for all $d\geq N$.
Likewise, the concentration-of-measure sequence
$1 - (1-\varepsilon^2/d)^{d/2}$ converges to $1 - e^{-\varepsilon^2/2}$
under the same Cauchy condition: an arbitrarily large fraction of
the uniform mass of $B^d$ sits within $\varepsilon/\sqrt{d}$ of the
boundary as $d$ grows.
\end{theorem}

\begin{proof}
$V_d = \pi^{d/2}/\Gamma(d/2+1)$ and $\Omega_{d-1} = 2\pi^{d/2}/\Gamma(d/2)$
are computable functions of $d$ (read $\pi$ and $\Gamma$ at the
relevant arguments via their partial-computation form,
Section~\ref{sec:integer-structure}). Stirling-type estimates give
computable bounds $V_d \leq C \cdot (2\pi e/d)^{d/2}$ for explicit
$C$, from which a computable $N(\varepsilon)$ such that
$V_d<\varepsilon$ for $d\geq N$ follows directly --- an
$\mathrm{RCA}_0$-internal statement. The same construction handles
$\Omega_{d-1}$ and the concentration-of-measure sequence.
\end{proof}

The convergence target $B^\infty$ has no volume to fill, no
boundary to inhabit, no interior to explore. It is the geometric
expression of ``distinction exists but content does not.'' It is
structured nothing --- the upper edge of the dimension tower, where
intrinsic content vanishes as $d$ grows.

% ============================================================
\section{Resolution: The Slicing Recurrence}\label{sec:nothing-unstable}
% ============================================================

The asymptotic upper edge $B^\infty$ is empty, but the dimension tower
as a whole is not inert. Each level has nontrivial intrinsic content
(volumes, sphere areas, beta-function ratios), and consecutive levels
are linked by the slicing recurrence. This recurrence is the
resolution mechanism: it relates the intrinsic content at $d+1$ to
the intrinsic content at $d$, extracting one dimension's worth of
finite content from the asymptotic limit at each step.

Cut $B^{d+1}$ perpendicular to one axis. Each cross-section at height
$x$ is a $d$-ball of radius $\sqrt{1-x^2}$. The volumes obey:
\[
V_{d+1} = V_d \cdot \sqrt{\pi}\cdot R(d+1), \qquad R(d) =
\frac{\Gamma((d+1)/2)}{\Gamma((d+2)/2)}.
\]
This is a sequence-level identity in the partial-computation form:
at any resolution depth $N$, the partial values $V_d^{(N)}$ obey
$V_{d+1}^{(N)} = V_d^{(N)}\cdot\sqrt{\pi}^{(N)}\cdot R(d+1)^{(N)} +
O(1/N)$, with the error bound $\mathrm{RCA}_0$-expressible from the
Riemann-sum approximations to the underlying Beta-function integral
$\int_{-1}^1 (1-x^2)^{d/2}\,dx$ and the partial half-integer Gamma.
The constant $\sqrt{\pi}$ is forced by orthogonality
(Theorem~\ref{thm:orth}): the angle between axis and equator is
$\pi/2$, forcing the half-integer argument in $B(1/2,\cdot)$, giving
$\Gamma(1/2) = \sqrt{\pi}$ as a convergence target.

Read structurally as the cascade's static content per dimension at any
fixed resolution stage: as $d$ ranges from $B^\infty$ down through
small finite values, volumes \emph{grow}, and the ball acquires content
it did not have at higher $d$. The volume peaks at $d = 5$. Each step
(from $d{+}1$ to $d$) extracts more finite content from the asymptotic
limit; the structured nothing of $B^\infty$ becomes structured
\emph{something} at every finite~$d$. The recurrence has no free
parameters: the resolution at each step is fully determined by the
Gamma function. (The temporal reading --- in which the universe's age
is the count of resolution stages reached so far, with truncation
height $N(t)$ growing one cascade layer per Planck tick --- is the
subject of the Remark below; the structural picture here is the static
content the dynamic process resolves into at each stage.)

\begin{remark}[Resolution as cosmic time]
The cascade's resolution process admits a physical reading developed
in Parts~II--VI. Under that reading, each step of the resolution ---
each new dimension added to the realised tower --- corresponds to
one Planck-time tick of cosmic evolution. The universe had a start
(when resolution began at finite stage), counts on indefinitely
(each tick adds one more dimension), and approaches but never reaches
$d=\infty$ (full resolution would be the asymptotic future, never
completed). The universe does not \emph{contain} infinity --- it is
the result of resolution-in-progress, observed from the current
stage. Time, on this reading, is not an additional ingredient: it
is the cascade's own resolution direction, registered as the
dimension tower's growth.

\smallskip
\noindent\textbf{Three vocabularies for one rate.} The downstream
papers develop the same resolution rate in three coordinated
readings, each offering a different way to picture what is happening:
\begin{itemize}
\item \emph{Mathematical (this Prelude, Part~0).} The slicing
recurrence (Section~\ref{sec:nothing-unstable}) extracts one
dimension's worth of finite content from $B^\infty$ at each step.
The recurrence is the resolution mechanism; the dimension tower
$\{(\mathbb{R}^d, B^d)\}_{d\in\mathbb{N}}$ is its trace; the integer
$d$ is the resolution stage. This is the structural reading: the
cascade as the unique parameter-free $d\to d-1$ recurrence on the
tower, indexed by Gamma-function values.
\item \emph{Physical (Parts~II and VI).} Each step is one Planck
tick of cosmic time. Part~II, Section~\xref{part2}{subsec:time-id} makes the static-versus-
dynamic identification explicit: each cascade slicing operation
\emph{is} one elementary unit of time --- ``slicing $=$ dimensional
collapse $=$ time'' --- not merely a coordinate parametrising the
tower. Part~VI promotes this to a cosmological dynamics: the
truncation height grows by one per tick,
$N(t+\alpha\,t_{\mathrm{Pl,red}}) = N(t)+1$, and the universe's
phases A--D are stages of the resolution prior to and beyond the
$\Omega_{217}$ floor. The universe's age is the count of resolution
stages reached so far.
\item \emph{Geometric (cover sheet).} The observer lives on the
$S^3$ shell of a 5D black hole, asymptotically falling toward the
pseudo-horizon. Each resolution step is one asymptotic step of that
infall: one dimension is collapsed at the cascade's boundary, and
the observer's host shell advances one Planck tick closer to the
pseudo-horizon. The infall never completes because $B^\infty$ has
infinitely many dimensions to resolve. ``Time dilation'' is the
local rate at which resolution proceeds along the observer's
worldline; ``the Big Bang'' is the resolution reaching the
$\Omega_{217}$ floor.
\end{itemize}
The three readings are dual descriptions of the same rate. Each
gives a distinct mental picture --- mathematical (a recurrence
extracting content), physical (a tick adding a layer), geometric
(an asymptotic infall) --- but they describe the same operation.

\smallskip
The Prelude commits only to the mathematical reading (the slicing
recurrence as the resolution mechanism, the dimension tower as its
trace); the physical identification of resolution with cosmic time,
the geometric identification with asymptotic infall toward the
pseudo-horizon, and the joint claim that one Planck time corresponds
to one tower extension, are the responsibility of the cover sheet
and Parts~II (time as dimensional collapse), V (cosmology, $\Omega_d$
at the cosmological-constant stage $d_2=217$), and~VI (tower growth,
phase structure of the pre- and post-Big-Bang universe).
\end{remark}

% ============================================================
\section{Where Resolution Lands: The Observer at \texorpdfstring{$d=4$}{d=4}}\label{sec:where-resolution-lands-the-observer-at-d4}
% ============================================================

The slicing recurrence resolves infinity into finite content at each
step. Where does the resolution stably terminate?

The volume maximum is at $d = 5$: the unique dimension where the unit
ball has the largest volume. The observer sits at $d = 4$, on the unit
sphere $S^3 = \partial B^4$ --- the boundary of the equatorial $B^4$-slice
of the volume-maximum ball $B^5$. By boundary dominance
($\Omega_{d-1}/V_d = d$ at $d=4$), this boundary carries
$\Omega_3/(\Omega_3+V_4)=4/5$ of the content.

Why does the observer end up here? Part~III of the series proves
that $d = 4$ is the unique spacetime dimension satisfying two
independent constraints simultaneously:

\begin{enumerate}
\item \textbf{Lovelock uniqueness}: the gravitational equation has a
unique form (Einstein's equation with cosmological constant) only at
$d = 4$.
\item \textbf{Complex spinors}: the Clifford algebra
$\mathrm{Cl}(1,d-1)$ has irreducibly complex minimal spinors --- the
prerequisite for quantum mechanics --- only when
$d\bmod 8\in\{4,5,6\}$.
\end{enumerate}

The intersection is $\{4\}$. No other dimension permits both unique
physical laws and quantum information processing. The observer is not
placed at $d = 4$ by assumption; $d = 4$ is where the resolution
trajectory naturally lands --- the unique stage at which stable
physical laws and complex quantum phase coincide. The resolution
process described in Section~\ref{sec:nothing-unstable} continues
beyond $d=4$ in both directions (toward $d=\infty$ as the asymptotic
upper edge, and toward $d=0$ as the trivial lower endpoint), but the
\emph{observable physics} of the resolution lives at $d=4$, the
boundary of the volume-maximum ball.

% ============================================================
\section{The Chain}\label{sec:the-chain}
% ============================================================

The complete logical chain, from pre-mathematics to physics, is:

\bigskip

\begin{center}
\renewcommand{\arraystretch}{1.4}
\resizebox{\textwidth}{!}{%
\begin{tabular}{rlll}
\toprule
Step & Statement & Status & Austerity \\
\midrule
0 & $0 \neq 1$ & Pre-mathematical & --- \\
1 & (B), (S), (L) properties of $\neq$; cumulation gives $d$-element $K_d$-shaped structure & Operational content of ``$\neq$'' as binary relation & --- \\
2 & Cardinality $d \in \mathbb{N}$ & Theorem~\ref{thm:iterate} & (i), (ii) \\
3 & Foundation: pure $\mathrm{RCA}_0$. All analytic content (sphere areas, ball volumes, $d\to\infty$ behaviour) is computable processes with $\mathrm{RCA}_0$-expressible convergence records & Strictly weaker than any classical analytic fragment, Section~\ref{sec:integer-structure} & (ii) \\
4 & $\pi$ forced as convergence target via Euler product on $\mathbb{N}$ ($\zeta_N(2),\,\mathcal{E}_N,\,\pi_N$ converge to a shared target) & Theorem~\ref{thm:pi-from-primes} (number theory on integers, no geometry, no completed real) & --- \\
5 & Vector-space realisation $\mathbb{R}^d$ & Hurwitz's theorem & (ii) \\
6 & Euclidean inner product on $\mathbb{R}^d$; angular extent $\pi$ inherited from step~4 & Elementary linear algebra & --- \\
7 & Original $d$ states at pairwise angle $\pi/2$ (extremal orthogonal) & Theorem~\ref{thm:orth} from (B), (S), (L) & --- \\
8 & Scale + rotation invariance at each $d\;\Rightarrow$ collapsed substrate; cascade as $d$-trajectory & Definition~\ref{def:ball}, Cor~\ref{cor:start} & (iii) \\
9 & Asymptotic $B^\infty$: zero volume, zero area & Theorem~\ref{thm:nothing} & --- \\
10 & Slicing recurrence, constant $\sqrt{\pi}$ & Beta-function on $\{B^d\}$ (steps~5--6) & --- \\
11 & Descent: volumes grow, peak at $d = 5$ & Gamma function & --- \\
12 & Observer at $d = 4$: boundary of the peak & Lovelock $\cap$ Clifford & --- \\
13 & Four distinguished dimensions: $5, 7, 19, 217$ & Part~0 & --- \\
14 & Cascade invariant $\approx 10^{-120}$ & Part~0 & --- \\
15 & Cosmological constant, QM, GR, SM, cosmology & Parts~I--V & --- \\
\bottomrule
\end{tabular}%
}
\end{center}

\bigskip

Steps 0--9 are established in this prelude. Step~10 is proved in
Part~0 (Theorem~3.1). Steps 11--12 are proved in Parts~0 and~III.
Steps 13--15 are the content of the series.

No step introduces a free parameter beyond the single declared
pre-mathematical input Definition~\ref{def:axiom} ($0\neq 1$), and no
step imports a foundational axiom beyond $\mathrm{RCA}_0$ itself.
Principle~\ref{princ:austerity} (austerity) is not a second input:
it is the operational content of Definition~\ref{def:axiom}'s status
as sole input, derived in Section~\ref{sec:what-nothing-means} above.
Steps invoking austerity clauses are noted in the ``Austerity''
column; the remaining chain steps are classical theorems on
integers, analysis, the Gamma function, or established uniqueness
results in known mathematics --- each read in the partial-computation
form of step~3, which is $\mathrm{RCA}_0$-internal. Beyond what
counting on $\mathbb{N}$ itself requires --- the arithmetic content
of $\mathrm{RCA}_0$ that cumulation produces --- the cascade therefore
imports \emph{no} further axiom: clause~(ii) of austerity rules out
the Monotone Convergence Theorem and stronger analytic principles,
because each analytic identity the cascade names is already an
$\mathrm{RCA}_0$-expressible statement about convergence of
computable sequences. $\mathrm{MCT}$, the
Riemann integral as a completed real, and related principles appear
only as bookkeeping shorthand in downstream papers (Parts~0--VI);
they are not foundational commitments (see status box and
Section~\ref{sec:open-questions}).

% ============================================================
\section{What This Paper Does Not Do}
\label{sec:what-not}
% ============================================================

\textbf{Does not derive physics.} This prelude contains no physical
predictions. Its purpose is to derive the starting structure of the
series --- the dimension tower
$\{(\mathbb{R}^d, B^d)\}_{d\in\mathbb{N}}$ of finite-dimensional
Euclidean spaces, with asymptotic upper edge $B^\infty$ --- from
$0\neq 1$ as the sole pre-mathematical input. Every step of the
chain is forced by clause~(ii) of austerity or by classical theorems
on integers and analysis; the chain table
(Section~\ref{sec:the-chain}) enumerates the specific clause invoked
at each step. In particular: clauses~(i) and~(ii) at step~2
(parameter economy forcing the dimension tower to extend without
bound; minimal strength forcing countable over uncountable indexing
via the $\mathrm{RCA}_0$ commitment); clause~(ii) at step~3
(selecting pure $\mathrm{RCA}_0$ as the foundational footprint:
every analytic identity the cascade names is an
$\mathrm{RCA}_0$-internal convergence statement about computable
sequences, so the Monotone Convergence Theorem and stronger
analytic axioms are not imported as foundational commitments,
Section~\ref{sec:integer-structure}); clause~(ii) at step~5
(selecting $\mathbb{R}$ over $\mathbb{C}$ or $\mathbb{H}$ via
Hurwitz's theorem); clause~(iii) at step~8 (unobservable quotient,
applied at each $d$ to both $\mathbb{R}^*$-dilation and
$\mathrm{SO}(d)$, collapsing the substrate at each $d$ to a single
point from which intrinsic content unfolds via the slicing
recurrence). Beyond what counting on $\mathbb{N}$ itself requires
--- the arithmetic content of $\mathrm{RCA}_0$ that cumulation
produces --- the cascade imports \emph{no} further axiom:
$\mathrm{MCT}$ and stronger analytic principles appear only as
bookkeeping shorthand in downstream papers (Parts~0--VI), not as
foundational commitments. The cascade series'
load-bearing hypothesis (cover sheet) is independent of the
prelude's chain: the downstream derivations take the dimension
tower (and its asymptotic upper edge $B^\infty$) as given and test
it by consequence.

\textbf{Real start, complex amplitudes, quaternionic spacetime.} The
Prelude's starting structure is over $\mathbb{R}$ at every level of
the dimension tower: no complex or quaternionic structure is
introduced by hand. By Hurwitz's theorem, the associative normed
division algebras over $\mathbb{R}$ are exactly $\mathbb{R}$,
$\mathbb{C}$, and $\mathbb{H}$, and the Gamma function governs the
slicing recurrence in all three cases. The cascade selects between
them downstream:
\begin{itemize}
\item \emph{Real at the start.} Each $\mathbb{R}^d$ in the dimension
tower is the most austere starting structure at dimension $d$ and is
the one used here. No complex or quaternionic structure is assumed at
any level.
\item \emph{Complex at the amplitude level.} A complex structure $J$
with $J^2=-\mathrm{Id}$ arises as a \emph{derived} consequence of the
forced precession $\alpha=\pi/2$ in Part~II. It is a theorem, not an
axiom: two consecutive cascade quarter-turns compose to multiplication
by $-1$, promoting the state space to a complex structure at the
quantum layer.
\item \emph{Quaternionic at the spacetime dimension.} Lovelock
uniqueness intersected with Clifford periodicity forces the spacetime
dimension to $d=4=\dim_\mathbb{R}\mathbb{H}$ in Part~III. The
\emph{spacetime} dimension matches the quaternions while the
\emph{amplitude} algebra remains $\mathbb{C}$.
\end{itemize}
The real/complex/quaternionic hierarchy is thus not an open question:
each algebra enters the cascade at the layer forced by a downstream
theorem. What the Prelude does not do is prove that the cascade could
not have been written over $\mathbb{C}$ or $\mathbb{H}$ from the
outset --- it proves only that the real case suffices, and that the
complex and quaternionic roles are recovered downstream without
additional input.

\textbf{Does not resolve the hard problem.} The cascade derives the
dimension where self-awareness is mathematically possible. It does
not derive self-awareness itself. The question ``why is there
something rather than nothing?''\ receives an answer: because
nothing, taken seriously, has structure. The question ``why is there
\emph{experience}?''\ remains open.

% ============================================================
\section{Open Questions}
\label{sec:open-questions}
% ============================================================

The prelude's chain from Definition~\ref{def:axiom} ($0\neq 1$) to
the cascade's complete starting structure
$\{(\mathbb{R}^d, B^d)\}_{d\in\mathbb{N}}$ with asymptotic upper edge
$B^\infty$ is closed at every step (Section~\ref{sec:the-chain}):
each step is forced by a clause of Principle~\ref{princ:austerity}
or by a classical theorem on integers and analysis. Beyond what
counting on $\mathbb{N}$ itself requires --- the arithmetic content
of $\mathrm{RCA}_0$ that cumulation produces --- the cascade imports
\emph{no} further axiom: clause~(ii) selects pure $\mathrm{RCA}_0$
as the foundational footprint, and every analytic identity the
cascade names is an $\mathrm{RCA}_0$-expressible convergence
statement about computable sequences --- no Monotone Convergence
Theorem, no Riemann integral as a completed real, no stronger
analytic principle is imported (Section~\ref{sec:integer-structure};
status box at the head of this prelude). $\mathrm{MCT}$ and its kin appear only as bookkeeping
shorthand in downstream papers (Parts~0--VI) when working through
Gamma identities, integrals, and asymptotic expansions; they are not
foundational commitments of the cascade.

\textbf{The universe as open-ended computation.} The partial-computation
reading dovetails with the cascade-is-resolution framing of
Section~\ref{sec:nothing-unstable}: the universe is in the middle of
resolving infinity, not at the asymptotic end. At any finite
resolution depth $N$ (the universe's current age in Planck ticks,
under the temporal reading) every cascade quantity --- $\pi$, $V_d$,
$\Omega_{d-1}$, the cascade invariant $10^{-120}$ --- has a specific
rational value $\pi_N$, $V_d^{(N)}$, etc., and the ``exact identities''
of standard mathematics are statements about the shared
convergence target as $N\to\infty$. The universe is not the
completed limit; it is the partial computation at the current $N$.

\textbf{The hard problem.} The cascade derives the dimension where
self-awareness is mathematically possible. It does not derive
self-awareness itself (already noted in
Section~\ref{sec:what-not}). The question ``why is there something
rather than nothing?''\ receives an answer: because nothing, taken
seriously, has structure. The question ``why is there
\emph{experience}?''\ remains open.

\bigskip

The series begins.

\begin{thebibliography}{99}

\bibitem{friedman1975} H.~Friedman, ``Some systems of second-order
arithmetic and their use,'' in \textit{Proceedings of the International
Congress of Mathematicians, Vancouver}, vol.~1, pp.~235--242 (1975).

\bibitem{simpson2009} S.~G.~Simpson, \textit{Subsystems of Second
Order Arithmetic}, 2nd ed., Cambridge University Press, 2009.

\end{thebibliography}

\end{document}
